% !TEX root = ../ts_notes.tex
\documentclass[12pt,fleqn]{article}
\usepackage{../vkCourseML}
\newcounter{chapter}
\hypersetup{unicode=true}
\interfootnotelinepenalty=10000
\newcommand{\dx}[1]{\,\mathrm{d}#1} % маленький отступ и прямая d
\begin{document}

Первая лекция состоит из двух частей.
Сначала обсуждаются организационные правила курса, затем вводится общий язык временных рядов: что считается рядом, какие задачи возникают, какие семейства моделей будут использоваться и зачем нужны сглаживание и декомпозиция.


\section{Организационные замечания}


В курсе предусмотрено пять домашних заданий.
Четыре из них будут практическими ноутбуками с задачами, одно - теоретическим набором задач.
Помимо домашних заданий будет контрольная работа.
Примеры прошлых контрольных и экзаменов доступны на Wiki курса; часть вариантов содержит решения.

Для домашних заданий действует политика просрочек: можно использовать либо две просрочки по 24 часа, либо одну просрочку на 48 часов.
Если контрольная пропущена по уважительной причине, ее вес можно перенести на экзамен, но для этого нужна стандартная справка и уведомление преподавателя.

Плагиат в домашних заданиях, контрольных и экзаменах строго наказывается.
Использование ChatGPT или других генеративных инструментов допустимо только как вспомогательный инструмент, например для генерации отдельных технических фрагментов кода.
Такие фрагменты нужно явно помечать и уметь объяснить.

Записи лекций и семинаров планируются, но их публикация зависит от посещаемости.
В конце домашних заданий будут формы обратной связи; преподаватель просит их заполнять, потому что они используются для улучшения курса.


\section{Определения}


В математическом анализе слово "ряд" обычно означает сумму.
В этом курсе под временным рядом понимается последовательность наблюдений во времени.

В общем виде случайный процесс можно записать как последовательность случайных величин

\[
Y_1,\ldots,Y_n
\]

или, если удобнее мыслить через индекс времени,

\[
(Y_t)_{t\in T}.
\]

Наблюдаемый временной ряд - это траектория, или реализация, этого случайного процесса:

\[
y_1,\ldots,y_n.
\]

Случайные величины принято обозначать заглавными буквами \(Y_t\), а конкретные наблюдения - строчными \(y_t\).
На практике часто приходится по одной траектории восстанавливать свойства всего процесса.
Это статистически наивная, но стандартная для временных рядов постановка: нескольких независимых траекторий одного и того же процесса обычно нет.

Курс в основном работает с низкочастотными данными: дневными, недельными, месячными, квартальными или годовыми.
Высокочастотные финансовые данные, например минутные котировки, почти не рассматриваются.
Для них нужны отдельные методы, ближе к pricing, stochastic calculus и quantitative finance.
В курсе финансовые данные появятся в основном в блоке про прогнозирование дисперсии и GARCH-модели.

Примеры низкочастотных рядов:

\begin{itemize}
  \item ВВП;
  \item инфляция;
  \item ключевая ставка;
  \item продажи товара по дням или неделям;
  \item выручка компании или банка;
  \item физические и научные измерения во времени.
\end{itemize}


\section{Возможные постановки задач}

\subsection{Случай одного ряда}


Самая очевидная задача - прогнозирование.
Если есть последовательность прошлых наблюдений, хочется предсказать будущие значения.
Но прогнозирование - не единственная задача анализа временных рядов.

Основные задачи для одного ряда:

\begin{itemize}
  \item прогнозирование будущих значений;
  \item заполнение пропусков;
  \item агрегация и дезагрегация;
  \item поиск выбросов и аномалий;
  \item фильтрация и декомпозиция.
\end{itemize}

Заполнение пропусков важно потому, что многие модели временных рядов требуют непрерывной временной сетки.
Как именно восстанавливать пропуски, зависит от структуры ряда: тренда, сезонности, частоты и наличия внешних факторов.

Агрегация переводит данные к более грубой частоте.
Например, часовые продажи можно просуммировать в дневные продажи.
Для переменных потока обычно берут сумму, а для переменных запаса чаще используют среднее или значение на конец периода.

Дезагрегация - обратная задача.
Например, из квартального ВВП нужно получить месячный ВВП.
Эта задача существенно сложнее: обычно приходится искать более частотные прокси-переменные и восстанавливать динамику по ним.

Декомпозиция раскладывает ряд на структурные компоненты: тренд, сезонность, цикличность и остаток.
Это важно не только для интерпретации, но и для прогноза: простые компоненты часто можно прогнозировать простыми моделями, а остаток затем моделировать отдельно или считать шумом.


\subsection{Случай набора рядов}


Если рядов много, можно решать те же задачи, что и для одного ряда, но появляются дополнительные постановки:

\begin{itemize}
  \item классификация рядов;
  \item кластеризация рядов;
  \item поиск похожих рядов;
  \item прогнозирование одного ряда с помощью других;
  \item многомерное моделирование;
  \item иерархическое прогнозирование.
\end{itemize}

Пример больших наборов рядов - соревнования Makridakis Competition, или M-Competition.
В M5 использовались иерархические данные Walmart: продажи множества товаров в разных магазинах, округах и штатах.
В таких задачах важно не только спрогнозировать отдельные ряды, но и согласовать прогнозы по уровням иерархии: прогнозы товаров должны суммироваться в прогнозы категорий, магазинов и регионов.

Многомерные модели, например VAR, позволяют моделировать несколько рядов совместно.
Обычно такие модели применимы к небольшому числу рядов; если рядов очень много, часто переходят к отдельным моделям, табличным методам или иерархическим схемам.


\section{Семейства моделей}


В курсе будут встречаться три больших семейства подходов.


\subsection{Классические статистические модели}


Классические модели рассматривают ряд как последовательность и напрямую описывают его временную структуру.
В курсе к ним относятся:

\begin{itemize}
  \item ETS;
  \item ARIMA и SARIMA;
  \item SARIMAX;
  \item GARCH и родственные модели.
\end{itemize}

Такие модели полезны, когда данных мало, нужна высокая скорость обучения или нужен надежный бенчмарк.
В типичных прикладных задачах низкочастотных рядов данных действительно немного: например, месячный ряд за несколько лет дает всего несколько десятков наблюдений.

ARIMA или auto-ARIMA часто используются как базовый ориентир даже тогда, когда итоговая модель будет сложнее.


\subsection{Табличное машинное обучение}


Второй подход - свести задачу временных рядов к табличной задаче.
Тогда можно использовать линейные модели, случайный лес, бустинги и другие стандартные методы.

В качестве признаков обычно берут:

\begin{itemize}
  \item лаги самого ряда;
  \item лаги внешних рядов;
  \item календарные признаки;
  \item признаки сезонности, например синусы и косинусы от времени;
  \item агрегаты по прошлым окнам.
\end{itemize}

Главная опасность такой постановки - утечка из будущего.
Признаки нужно строить так, чтобы в момент прогноза модель не использовала информацию, которая в реальности еще недоступна.

Табличные модели часто оказываются очень сильным и гибким инструментом, особенно если классические статистические модели слишком ограничены.


\subsection{Deep Learning}


К нейросетевым моделям для временных рядов относятся RNN, LSTM и Transformer-подходы.
Они могут быть полезны, когда данных очень много или частота наблюдений высокая.
Но для стандартных низкочастотных рядов такие модели часто неудобно обучать и настраивать, а выигрыш над бустингами или классическими моделями не гарантирован.

Поэтому основной фокус курса будет на классических статистических моделях и табличном машинном обучении.


\section{Периодичность и частота}


Большинство моделей временных рядов предполагает, что наблюдения идут через фиксированные промежутки времени.
Такой ряд называется периодическим, или регулярным.
Примеры:

\begin{itemize}
  \item каждый день;
  \item каждую неделю;
  \item каждый месяц;
  \item каждый квартал;
  \item каждый год.
\end{itemize}

Если промежутки между наблюдениями нерегулярны, ряд можно назвать непериодическим или спорадическим.
Для таких данных есть отдельные модели, но в прикладных экономических задачах чаще пытаются привести данные к регулярной сетке.

Частота - это период, с которым появляются наблюдения.
У годовых данных естественная частота одна: год.
У дневных данных одновременно могут проявляться несколько сезонных частот: недельная, месячная, квартальная и годовая.
Это важно для последующей декомпозиции и для алгоритмов вроде STL и MSTL.
\newpage
% \begin{thebibliography}{0}
% \end{thebibliography}

\end{document}
